\section{Local control without nondegeneracy}
\label{sec:local}

We establish uniform bounds for the local groups of all iterates of a
finite collection of prime geodesics.  The bounds include geodesics of
zero mean index.  We also retain the actual restriction maps to the
half-translation fixed locus and show that evaluation cup product
vanishes at each reflection crossing.  The analytic preliminaries below
apply to every smooth Riemannian metric, independently of the
finite-geodesic hypothesis.

\subsection{Energy sublevels and finite-dimensional models}

\begin{lemma}[Analytic foundations]\label{lem:local-analytic}
The energy \(E\) is smooth on \(\Lambda=H^1(S^1,S^n)\), satisfies the
Palais--Smale condition, and has precisely the smooth closed geodesics
and the constants as critical points.  There is a complete strong
Riemannian metric on \(\Lambda\), invariant under \(G=O(2)\), for which
a bounded normalized negative gradient of \(E\) has a complete
equivariant flow.  Sufficiently small energy sublevels retract
equivariantly to \(M\).
\end{lemma}

\begin{proof}
Embed the sphere in \(\mathbb R^{n+1}\).  Pointwise exponential charts,
the embedding \(H^1(S^1)\subset C^0(S^1)\), and the one-dimensional
\(H^1\) product rule make \(\Lambda\) a smooth Hilbert submanifold of
\(H^1(S^1,\mathbb R^{n+1})\).  Its induced strong metric is complete.
Indeed, an intrinsic Cauchy sequence is ambient Cauchy; its limit still
takes values in the closed sphere, and a local Hilbert chart near that
limit bounds intrinsic path distances by ambient distances.  Time
translations and reflection are isometries.  Each acts smoothly, and
their action is jointly continuous in the strong \(H^1\) topology.

Extend \(g\) to a smooth positive symmetric matrix \(A(x)\) along the
ambient tangent bundle, preserving the tangent--normal splitting.
Then
\[
 E(\gamma)=\int_0^1
       \langle A(\gamma)\dot\gamma,\dot\gamma\rangle\,\dd t
\]
is smooth.  To check compactness, suppose that \(E(\gamma_r)\) is bounded
and \(\| \dd E(\gamma_r)\|\to0\).  Pass to a subsequence converging weakly
in \(H^1\) and uniformly to \(\gamma\).  Let \(P_x\) denote Euclidean
tangent projection.  The tangent fields
\[
 \eta_r=P_{\gamma_r}(\gamma_r-\gamma)
\]
are bounded in \(H^1\).  Testing \(\dd E(\gamma_r)\) on \(\eta_r\), every
term containing \(\gamma_r-\gamma\) without a derivative tends to zero:
its uniform norm tends to zero, while the squared velocities have
bounded integrals.  This includes the terms from differentiating
\(A(\gamma_r)\) and \(P_{\gamma_r}\).  Consequently
\[
 \int_0^1\!
 \langle A(\gamma_r)\dot\gamma_r,
                 \dot\gamma_r-\dot\gamma\rangle\,\dd t\longrightarrow0.
\]
Uniform convergence of the coefficients and weak convergence of the
velocities now give convergence of their energy norms.  Uniform
positivity of \(A\) proves strong \(H^1\) convergence.

At a critical point the weak Euler--Lagrange equation says, in a
coordinate interval, that the distributional derivative of the metric
momentum is an \(L^1\) quadratic expression in the velocity.  The
momentum, hence the velocity, is continuous.  The geodesic equation
then gives smoothness.  Its periodic boundary condition follows from
the first variation.  Conversely, geodesics and constants are critical.

Normalize the negative gradient by
\[
 Z=-\frac{\nabla E}{\sqrt{1+\|\nabla E\|^2}}.
\]
This is a smooth invariant vector field of norm at most one, strict
away from the critical set.  A bounded vector field on a complete
strong Riemannian manifold has a complete flow: a trajectory of finite
time has finite length and a limit, through which the local flow
extends it.

Finally, near the constants use coordinates
\(\gamma(t)=\exp_x \xi(t)\), where \(\xi\in H^1(S^1,T_xM)\) has mean
zero.  The inverse function theorem gives a uniform tube over the
compact zero section.  The normal Hessian of energy is positive
coercive by the mean-zero Poincar\'e inequality.  Taylor integration
and the positive square-root change of normal coordinates make the
energy an exact positive quadratic form there.  These constructions
are equivariant.  Small energy forces both the diameter and the
ambient \(L^2\) derivative norm of a loop to be small, so all sufficiently
small sublevels lie in this tube.  Radial contraction gives the claim.
\end{proof}

\begin{lemma}[Regular strips]\label{lem:local-strips}
If a closed bounded energy interval contains no critical value, the
gradient flow gives equivariant sublevel equivalences across it.  At
regular cuts, open and closed sublevel pairs compute the same
equivariant cohomology, naturally for inclusions and fixed loci.
\end{lemma}

\begin{proof}
Palais--Smale implies a positive lower bound on \(\|\nabla E\|\) in
such an interval.  The normalized flow therefore decreases energy at
a uniformly positive rate there.  A fixed sufficiently long time
moves its upper sublevel into its lower sublevel, and the intervening
flow gives the required homotopies.  A small regular collar at an
endpoint similarly compares the open and closed sublevels.  Apply
these constructions simultaneously at the endpoints of a pair.
All maps are equivariant because they use the same invariant flow.
\end{proof}

We shall use closed regular polygon sublevels in the local
constructions.  Lemma~\ref{lem:local-strips} identifies their crossing
groups with those of the open sublevels \(X_a=\{F<a\}\) used elsewhere.
Here and below, a cut for polygon energy is the square of the
corresponding cut for \(F=\sqrt E\).

\begin{lemma}[Polygon models]\label{lem:local-polygons}
Fix \(B<\infty\).  For every sufficiently large integer \(N\), the
energy sublevel \(\{E\le B^2\}\) retracts onto its broken-geodesic
polygons with \(N\) equal time cells.  The retraction preserves every
smaller energy sublevel.  It commutes with reflection and all time
shifts preserving the vertex grid, hence with \(K\) when \(N\) is even.
Regular polygon sublevels through this bound are compact smooth
manifolds with boundary.
At every critical point in the band, polygon and Hilbert indices and
nullities agree.  Regular energy cuts are dense in the band, and
their sublevels have finite CW models.
\end{lemma}

\begin{proof}
Choose \(\rho>0\) below the injectivity radius.  The open manifold
\[
 P_N=\{(x_0,\ldots,x_{N-1})\in M^N:
                       d_g(x_j,x_{j+1})<\rho\}
\]
is identified with the polygons formed by minimizing constant-speed
segments; subscripts are cyclic.  Its energy is
\[
 f_N=N\sum_{j=0}^{N-1}d_g(x_j,x_{j+1})^2.
\]
Choose \(N\) so large that \(B/\sqrt N<\rho\).  On each time cell,
replace the central subinterval of proportion \(s\), \(0\le s\le1\),
by its minimizing segment.  Cauchy--Schwarz makes all endpoint
distances shorter than \(\rho\), and shows that this replacement does
not increase energy.  At \(s=1\) the loop is a polygon, and polygons
are fixed throughout.

For positive subinterval width, strong \(H^1\) continuity follows from
smooth dependence of short geodesics on endpoints and continuity of
restriction and rescaling in \(L^2\).  As the width tends to zero, the
replacement energy is bounded by the original energy on the
disappearing intervals.  Those integrals tend uniformly to zero along
any strongly \(H^1\)-convergent sequence.  This proves continuity also
at \(s=0\).  Centering the replacements proves the stated symmetries.

On the closed sublevel \(f_N\le B^2\), every edge has length at most
\(B/\sqrt N<\rho\), so this sublevel is compactly contained in \(P_N\).
Its boundary is smooth at a regular value.  The polygon first
variation imposes matching velocities at every vertex, so its
critical points are exactly the smooth closed geodesics in the band.

At such a geodesic, decompose an \(H^1\) variation field into its
piecewise Jacobi interpolation and a field vanishing at every
vertex.  Integration by parts makes the two summands orthogonal for
the index form.  On a cell the second summand satisfies the
Dirichlet Poincar\'e inequality.  A curvature upper bound, the speed
bound \(B\), and sufficiently small cell width make its index form
positive definite.  The polygon Hessian is the first summand, which
therefore has the full Hilbert index and nullity.
Sard's theorem for \(f_N\) gives dense regular cuts, since its critical
values in the band are exactly those of \(E\).  Compact smooth
manifolds with boundary have finite CW models, completing the claim.
\end{proof}

If \(k\mid N\), the fixed manifold of shift by \(1/k\) in \(P_N\)
is the repeated copy of \(P_{N/k}\), and
\begin{equation}\label{eq:local-energy-scaling}
 f_N|_{P_{N/k}}=k^2f_{N/k}.
\end{equation}
In particular, the half-translation fixed manifold has energy
\(4f_{N/2}\).  These identities follow directly from the sum defining
\(f_N\).

\subsection{Prime images and iteration}

\begin{lemma}[Counting prime images]\label{lem:local-prime-images}
On any smooth Riemannian manifold, every nonconstant closed geodesic
is an iterate of a prime geodesic.
Two prime geodesics with the same image differ by phase and possibly
reversal.  Reversal never preserves an oriented phase orbit.
\end{lemma}

\begin{proof}
A nonconstant geodesic has constant positive speed and is an
immersion.  Its periods form a closed discrete subgroup of
\(\mathbb R\), so there is a least positive period.  Dividing by that
period gives its prime root.

To compare two primes with the same image, first use unit-speed
parameters.  Cover the parameter circle of the first by finitely many
compact intervals on which it is an embedded arc.  The inverse
images of those arcs under the second geodesic form a finite closed
cover of its parameter circle.  By Baire's theorem, one contains an
interval.  Within this interval choose a point whose image is not an
endpoint of the chosen arc; an immersion cannot take an interval
into that finite endpoint set.  Near this point both geodesics lie
in the same embedded arc and have the same tangent line.  Their
unit velocities agree up to sign.  Uniqueness for the geodesic
equation identifies them globally up to translation and reversal.
Their least unit-speed periods agree, giving the assertion for
period-one parametrizations as well.  This proof allows
self-intersections of either prime.

If reversal preserved an oriented phase orbit, one would have
\(\gamma(t)=\gamma(a-t)\).  Differentiation at a fixed point of this
parameter reflection gives \(\dot\gamma=-\dot\gamma\), impossible for
a nonconstant geodesic.
\end{proof}

Suppose henceforth, when discussing a finite collection of critical
labels, that there are only finitely many prime images.  Let
\(c_1,\ldots,c_p\) be prime representatives and \(\ell_i>0\) their
lengths.  By Lemma~\ref{lem:local-prime-images}, the critical set at
positive energy consists of the two oriented phase circles for
each iterate \(c_i^m\).  Its critical \(F\)-values are \(m\ell_i\).
There are at most \(2p\) oriented labels at a level and
\begin{equation}\label{eq:local-event-count}
 O(1+b)\quad\hbox{labels below \(b\)},\qquad
 O(1+D)\quad\hbox{labels in an interval of width \(D\)}.
\end{equation}
The constants depend only on the finite list of lengths.  In
particular, positive critical values are locally finite.

Write \(I_{i,m}\) for the periodic Morse index of \(c_i^m\).
The following elementary form of the iteration estimate suffices;
it is consistent with the sharper classical iteration theory
\cite{Bott1956Iteration}.

\begin{lemma}[Index and root types]\label{lem:local-iteration}
For each prime \(c_i\), there is a number \(\Delta_i\ge0\) such that
\begin{equation}\label{eq:local-index}
 |I_{i,m+k}-I_{i,m}-I_{i,k}|\le n,\qquad
 |I_{i,m}-m\Delta_i|\le n.
\end{equation}
There is an integer \(D_i\ge1\) such that, for
\(s=\gcd(m,D_i)\), every null field at \(c_i^m\) repeats a null field
at \(c_i^s\).  Thus only finitely many root kernels occur.
For even \(m\),
\begin{equation}\label{eq:local-even-index}
 |I_{i,m}-2I_{i,m/2}|\le n.
\end{equation}
\end{lemma}

\begin{proof}
Use the periodic index form over integer intervals at prime speed.
In the form domain for the \((m+k)\)-period problem, require that the
value at time \(m\) equal the initial value.  In the direct sum of
the \(m\)- and \(k\)-period domains, require that the two initial
values agree.  Both are restrictions of codimension at most \(n\),
and the restricted forms identify.  If the two indices drop by
\(a,b\in[0,n]\), their original difference is \(a-b\).  This proves
the first estimate.  The indices are finite since the index form
is a positive derivative form plus a compact zeroth-order
perturbation.

The sequences \(I_{i,m}+n\) and \(I_{i,m}-n\) are respectively
subadditive and superadditive.  Their normalized limits agree;
applying the inequalities along multiples of each fixed \(m\)
gives
\[
 \frac{I_{i,m}-n}{m}\le\Delta_i
                         \le\frac{I_{i,m}+n}{m}.
\]
Nonnegative indices give \(\Delta_i\ge0\).
Equation~\eqref{eq:local-even-index} is the first estimate with two
equal summands.

Let \(P_i\) be the real Jacobi Cauchy-data monodromy over one prime
period.  The periodic kernel is \(\ker(P_i^m-1)\).  Take \(D_i\) to
be the least common multiple of the orders of its root-of-unity
eigenvalues, with value one if there are none.  For
\(s=\gcd(m,D_i)\), the same eigenvalues satisfy \(\lambda^m=1\)
and \(\lambda^s=1\).  On a contributing complex Jordan block,
\(z^m-1\) has a simple root, so its kernel is the ordinary
eigenspace.  All other blocks contribute zero.  Hence
\[
 \ker(P_i^m-1)=\ker(P_i^s-1).
\]
ODE uniqueness identifies the corresponding fields as repeated
root fields.  This computation is over \(\mathbb R\) or
\(\mathbb C\); it makes no division by multiplicity in
\(\kk\)-cohomology.
\end{proof}

In particular, \(\Delta_i=0\) implies \(0\le I_{i,m}\le n\) for
every \(m\).  We shall retain these iterates; they may occur at
arbitrarily large lengths.

\subsection{A common flow for local restriction maps}

The next lemma is also useful for a compact Morse--Bott critical
manifold.  It does not require finitely many prime geodesics.
Here a \emph{good pair} means a closed pair whose second space
is a strong deformation retract of an open neighborhood; its
relative singular cohomology is the reduced cohomology of its
pointed quotient.

\begin{lemma}[Flow lenses]\label{lem:local-lens}
Let a compact group \(H\) act on a finite-dimensional or Hilbert
manifold, and let \(f\) be a smooth \(H\)-invariant function with a
complete bounded \(H\)-equivariant strict descending flow \(\phi\).
Suppose the critical set
in a closed energy band is a finite disjoint union of compact
sets at a single level \(c\), and that \(H\) permutes these sets.
Assume that the energy drop is
bounded below away from every neighborhood of those sets, and
that the corresponding assertion holds on the annuli of nested
neighborhoods.  These conditions hold for the normalized energy
flow of Lemma~\ref{lem:local-analytic}, and on compact polygon
bands.

For a common sufficiently large \(T\) and sufficiently tight
regular cuts \(a<c<b\), set
\[
 X=\{f\le b\},\quad X_-=\{f\le a\},\quad f_T=f\circ\phi_T,
\]
\begin{equation}\label{eq:local-lens}
 Y=\{x\in X:f_T(x)\le a\},\qquad
 N=\{x\in X:f_T(x)\ge a\},\qquad Q=N\cap Y.
\end{equation}
There are natural relative cohomology isomorphisms, also after
Borel construction, induced by
\begin{equation}\label{eq:local-lens-maps}
 (X,X_-)\longrightarrow(X,Y)\longleftarrow(N,Q).
\end{equation}
The lens splits into clopen pairs in the chosen disjoint critical
neighborhoods.  All maps commute with restriction to a fixed
submanifold when the same flow is restricted there, with subgroup
restriction, and with cup product by ambient absolute classes.
\end{lemma}

\begin{proof}
Choose disjoint neighborhoods \(W_\alpha\) of the compact critical
sets, and smaller \(V_\alpha\) with positive distance from the
complement of \(W_\alpha\); the choices are equivariantly
permuted.  Work first in a fixed isolating band.  Uniform drop
outside the \(V_\alpha\) makes a sufficiently long segment
remaining in that band visit one of them.  Traversing the annulus
between \(V_\alpha\) and the complement of \(W_\alpha\) takes
positive time, since the speed is bounded and its width is
positive.  The traversal costs a uniformly positive amount of
energy.  Choose the final width \(b-a\) smaller than all these
costs.  Every segment of duration \(T\) with initial energy at
most \(b\) and terminal energy at least \(a\) is then trapped
in a single \(W_\alpha\), before and after its visit to
\(V_\alpha\).

The first map in \eqref{eq:local-lens-maps} is a pair homotopy
equivalence: its inverse is the map induced by \(\phi_T\).
The intermediate flow preserves the relevant subspaces.
The finite closed cover \(X=N\cup Y\) gives
\[
 N/Q\cong X/Y
\]
by closed pasting.  This does not require \(X\) to be compact.
At \(f_T=a\), the derivative of \(f_T\) along the flow is
strictly negative.  Transverse first hitting of this level is
continuous, with hitting time zero on the lower side.  Flowing
for a fraction of that time gives a neighborhood retraction
onto \(Y\) in \(X\).  On the upper side, stop at equality to
obtain the corresponding retraction onto \(Q\) in \(N\).
The condition \(f\le b\) is preserved by descent.  These are
good pairs, so their quotient homeomorphism gives the second
relative isomorphism.  If a collapse set is empty, adjoin a
disjoint basepoint.

The constructions are equivariant.  Finite invariant closed
covers remain closed covers after Borel construction, and the
neighborhood homotopies descend.  This proves the equivariant
version without asserting compactness of a Borel space.
Trapping makes \(N\cap W_\alpha\) and \(Q\cap W_\alpha\) clopen
label pairs.  Restriction to a fixed manifold intersects every
space in \eqref{eq:local-lens-maps} with that manifold, using
the original flow.  Finally, the maps in that diagram are
inclusions and hence respect ambient cup products and
connecting homomorphisms.
\end{proof}

For the normalized energy flow, the drop hypotheses in
Lemma~\ref{lem:local-lens} follow from Palais--Smale: a sequence
away from the chosen critical neighborhoods with drop tending
to zero would converge to a critical point outside them.
For a finite polygon band they follow from compactness.

\begin{theorem}[Uniform local control]\label{thm:local-control}
Assume there are finitely many prime images.  There are constants
\(A,B<\infty\) such that every sufficiently tight crossing of a
positive critical level has the following properties.
\begin{enumerate}[label=\textup{(\roman*)}]
\item It decomposes into unoriented labels \((i,m)\) at that
level.  Let \(\mathcal P_{i,m}\) be the ordinary pair for one
orientation.  Then
\[
 \sum_j\dim H^j(\mathcal P_{i,m};\kk)\le B,\qquad
 H^j(\mathcal P_{i,m};\kk)=0
                  \quad\text{if }|j-I_{i,m}|>A.
\]
\item Its reflection block is induced from \(\mathcal P_{i,m}\);
its \(K\)-block is induced from the
\(\langle\tau\rangle\)-action on that pair.  For odd \(m\)
the fixed label is empty.  For even \(m\), the ordinary fixed
pair has total rank at most \(B\) and support
\(|j-I_{i,m/2}|\le A\).
\item Full-to-fixed restriction respects the labels.  These
identifications are compatible with ambient maps and
connecting homomorphisms.
\item Cup product by
\(e=\ev_0^*[S^n]^*\in H_R^n(\Lambda;\kk)\) vanishes on
the reflection crossing, including each individual label.
\end{enumerate}
The constants depend on the metric and the finite prime list,
and include all degenerate iterates and all zero-mean-index
primes.
\end{theorem}

\begin{proof}[Proof of the decomposition, naturality, and cup assertion]
Choose an even polygon order adequate for the event.  Start with
one complete \(K\)-invariant strict descending gradient flow
for the original \(f_N\).  Its gradient is tangent to
\(\operatorname{Fix}(\tau)\).  Thus a critical point of the
restriction is exactly a full critical point there, and the
restricted flow is strict away from them.

By Lemma~\ref{lem:local-prime-images}, the critical circles
come in disjoint reversed pairs.  The half-shift preserves
each orientation circle and fixes it pointwise exactly when
\(m\) is even; for odd \(m\), choose its neighborhood disjoint
from the fixed locus.  Use Lemma~\ref{lem:local-lens} with
disjoint equivariantly permuted neighborhoods of these
oriented circles.  Reflection induction identifies a block
with the ordinary pair of one orientation, while \(K\)
induction retains its half-shift action.  Literal intersection
of the original lenses with the fixed manifold gives the
full-to-fixed map.  In particular, the actual diagrams are
\[
\begin{tikzcd}
 H_K^j(X,X_-) \arrow[r,"\cong"] \arrow[d,"\res"']
   & H_K^j(N,Q) \arrow[d,"\res"]\\
 H_K^j(X^\tau,X_-^\tau) \arrow[r,"\cong"]
   & H_K^j(N^\tau,Q^\tau),
\end{tikzcd}
\]
with \(N^\tau=N\cap\operatorname{Fix}(\tau)\) and likewise
for \(Q\).  They commute with ambient relative maps and
the associated exact sequences.  No alternative local
normal form is being used to identify these maps.

The finite union of prime images does not fill \(S^n\), for
\(n\ge2\).  Choose \(q\) outside it, and choose all critical
neighborhoods inside the invariant open set
\[
 \bigcap_{g\in K}g^{-1}
       \bigl(\ev_0^{-1}(S^n\setminus\{q\})\bigr).
\]
The entire lens therefore evaluates into the punctured
sphere.  Evaluation at zero is reflection invariant, so it
induces \(ER\times_R N\to S^n\setminus\{q\}\).  The top
sphere class pulls back to zero.  Hence its cup product on
\(H_R^*(N,Q)\) is zero.  The inclusion isomorphisms of
Lemma~\ref{lem:local-lens} transfer this vanishing to the
actual crossing.  This uses only reflection invariance of
evaluation; half-translation invariance is not needed.
All choices are common to the finitely many labels at a
simultaneous level.  The uniform rank statements are proved
next.
\end{proof}

\subsection{Computing ranks while retaining the original maps}

The finite-root reduction is the local mechanism in classical
iteration arguments \cite{GromollMeyer1969}.  We give the required
rank calculation explicitly, including its relative Thom pair.

Fix one oriented label at energy \(c\).  Add small nonnegative
bumps, constant near the other circles at \(c\), supported away
from the selected neighborhood \(W\).  On their compact
transition regions the original strict drop is bounded away
from zero.  Taking the amplitudes sufficiently small leaves
the same original flow strict for the modified energy \(f_i\).
No new critical points appear, and only the selected circle
remains in a sufficiently narrow band about \(c\).

Choose one large time and tight cuts for the original energy
and this finite list of bumped energies.  The total lens for
\(f_i\) is exactly its original labeled lens: in either
direction every relevant segment is trapped in \(W\), where
the functions agree.  The same construction on the fixed
manifold identifies its ordinary label with an isolated
crossing.  Its ordinary rank is now independent of the strict
flow and of smaller regular cuts, by strip deformations.
We may compute these ranks using a split flow.  The original
flow remains the one defining all maps in
Theorem~\ref{thm:local-control}.

\begin{lemma}[The normal shift]\label{lem:local-shift}
Put \(s=\gcd(m,D_i)\) and \(k=m/s\).  The ordinary local rank
at \(c_i^m\), in degree \(j\), equals the ordinary local rank
of the corresponding root label at \(c_i^s\), in degree
\(j-(I_{i,m}-I_{i,s})\).
\end{lemma}

\begin{proof}
\emph{Normal coordinates.}
Take \(N\) divisible by \(k\) and sufficiently large.
The fixed manifold \(Z=P_{N/k}\subset P_N\) of the
\(1/k\)-shift has restricted energy
\(h=k^2f_{N/k}\).  Near the selected orbit the bump vanishes.
Use a cyclically invariant tubular metric.  Symmetry makes
the normal first derivative vanish along \(Z\).  The
Hessian has reducing fixed and normal summands, and
Lemma~\ref{lem:local-iteration} places its full kernel
inside the fixed tangent space.  The normal block is
therefore nonsingular, of negative rank
\[
 b_-=I_{i,m}-I_{i,s}\ge0.
\]

We record the normal coordinate construction to keep the
degenerate base unchanged.  Taylor integration writes
\[
 f_i(z,\nu)=h(z)+\langle A(z,\nu)\nu,\nu\rangle,
\]
where \(A_0(z)=A(z,0)\) is symmetric and invertible.
Normalize by \(|A_0|^{1/2}\), putting
\[
 B=|A_0|^{-1/2}A|A_0|^{-1/2},\qquad
 J=\operatorname{sign}(A_0).
\]
Near the zero section, \(B\) is close to \(J\).  The
near-identity series \(D=(JB)^{1/2}\) is \(J\)-self-adjoint,
and \(D^*JD=B\).  Thus the fiber inverse function theorem
applied to \(D|A_0|^{1/2}\nu\) gives
\begin{equation}\label{eq:local-split}
 f_i=h(z)-|u|^2+|y|^2.
\end{equation}
The construction is uniform near the compact orbit.
The positive and negative bundles can be nontrivial; no
trivialization of them is required.

\emph{The local quotient.}
Let \(\psi_t\) be a complete strict descending flow for the
unmodified \(h\).  In \eqref{eq:local-split}, lift it using
metric connections, expand \(u\) exponentially, and contract
\(y\) exponentially.  The energy derivative is
\[
 \dd h(\dot\psi)-2|u|^2-2|y|^2<0
\]
off the critical set.  Convex interpolation with a
normalized negative gradient gives a complete strict
full field agreeing with this one near the orbit.

Choose a full trap \(W_f\) whose base projection has compact
closure inside the base chart and avoids all other
\(h\)-critical points.  The exact-field region contains
uniform fiber balls over this closure.  A smaller selected
base trap \(W_h\) lies under uniform small fiber balls
contained in \(W_f\); all other base traps lie outside
its projection.  First choose a common sufficiently large
time \(T\) for full and base trapping, then sufficiently
tight regular cuts \(a<c<b\).  In particular,
\(\sqrt{b-a}\) is smaller than the available fiber radius.

Apply Lemma~\ref{lem:local-lens} to the isolated full
crossing.  On its quotient \(N/Q=X/Y\), shrink \(y\)
to zero.  The new first \(T\)-paths have the same base
and negative coordinate, and a scaled positive coordinate,
so they stay in the exact region.  Both initial and terminal
energies decrease.  Thus the maps take \(N\) into \(X\)
and \(Q\) into \(Y\), and descend continuously to the
quotient; some points may collapse.  This retracts onto
the \(y=0\) subset together with the basepoint.

Apply the full forward flow to this subset.  While a
point remains uncollapsed, every shifted \(T\)-segment
is trapped in the exact region, so \(y=0\) persists and
the base follows \(\psi\).  On the compact base projection
where \(h\ge b>c\), its drop has a positive lower bound.
A single further time therefore sends all survivors
into \(h\le b\).  This last subset is preserved up to
collapse, so its inclusion is a homotopy equivalence.
The homotopies are continuous on the pointed quotient;
products of a quotient map with the compact interval
remain quotient maps.

Write \(h_T=h\circ\psi_T\).  The invariant subset \(h\le b\)
used in this homotopy equivalence, with its equality boundary before collapse,
is described by
\begin{equation}\label{eq:local-base-lens}
 P=\{h\le b,\ h_T\ge a\}_{\mathrm{selected}},\qquad
 A_P=\{h_T=a\}\cap P,\qquad
 |u|\le e^{-T}\sqrt{h_T(z)-a}.
\end{equation}
For one direction, the projected segment stays in the
base band and the projection of the full trap, so it
belongs to the selected base label.  Conversely, a
selected base path is trapped in \(W_h\), and
\(e^t|u|\le\sqrt{b-a}\) for \(0\le t\le T\).
Its lift therefore stays in the exact region and
satisfies the full lens inequalities.  Equality in
the radius bound is precisely its exit.

Let \(D_-\) and \(S_-\) be the unit disk and sphere
bundles of the negative bundle over \(P\).  Scaling by
the continuous radius in \eqref{eq:local-base-lens}
gives a homeomorphism from
\begin{equation}\label{eq:local-thom-quotient}
 D_-/(S_-\cup(D_-)_{A_P})
\end{equation}
onto the final pointed subset.  Indeed it is bijective
away from the collapsed union; its domain is compact
and its target is Hausdorff.  The radius vanishes
exactly over \(A_P\).  This also handles rank zero
and the disjoint-basepoint convention.

\emph{Relative Thom calculation.}
For completeness, we prove the relative Thom calculation
for this lens, without requiring a manifold structure on
its base.  On \(P\), let \(t_p\) be the first forward
time until \(h_T\le a\), truncated at a fixed \(T_0>0\).
It is continuous.  Finite hits are transverse since
\(a\) is regular; at a point that has not hit by
\(T_0\), either the hit occurs transversely at \(T_0\),
or the strict terminal inequality persists in a
neighborhood.  The base flow stays in its clopen
label until equality.

For a unit-disk coordinate put
\[
 t(z,u)=\min\{t_p(z),-\log|u|\},\qquad -\log0=+\infty.
\]
This is continuous also at \(u=0\), where the radial
entry is inactive for \(|u|<e^{-T_0}\).
Its zero set is
\(\mathcal A=S_-\cup(D_-)_{A_P}\).
On the open neighborhood \(t<T_0\), flow the base
for time \(\lambda t\), parallel transport \(u\),
and multiply its norm by \(e^{\lambda t}\).
Neither the base exit nor the sphere is passed before
the stopping time.  The residual minimum time is
\((1-\lambda)t\), so the neighborhood is preserved,
the union is fixed, and at \(\lambda=1\) it is reached.
This proves that \((D_-,\mathcal A)\) is a good pair.
Rank zero simply removes the radial alternative.

The sphere pairs are good by radial collars.
Closed pasting identifies
\[
 H^*(\mathcal A,S_-)
       \cong H^*((D_-)_{A_P},(S_-)_{A_P}).
\]
Under the natural absolute mod-2 Thom isomorphisms,
restriction from \(H^*(D_-,S_-)\) to this group is
the restriction \(H^{*-b_-}(P)\to H^{*-b_-}(A_P)\).
The exact sequence of the triple
\(S_-\subset\mathcal A\subset D_-\) gives
\[
\begin{split}
0\longrightarrow&
\coker\bigl(H^{j-b_--1}(P)\to H^{j-b_--1}(A_P)\bigr)
\longrightarrow H^j(D_-,\mathcal A)\\
\longrightarrow&
\ker\bigl(H^{j-b_-}(P)\to H^{j-b_-}(A_P)\bigr)
\longrightarrow0.
\end{split}
\]
The base-pair exact sequence has the same kernel
and cokernel.  Thus
\begin{equation}\label{eq:local-thom-dimension}
 \dim H^j(D_-,\mathcal A)
                    =\dim H^{j-b_-}(P,A_P).
\end{equation}
This argument compares vector spaces, rather than
subtracting possibly infinite dimensions.

Absolute Thom applies to these compact metrizable
bases as follows.  Pull the bundle back along the
CW weak equivalence \(|\operatorname{Sing}P|\to P\).
The disk and sphere bundle projections are fibrations,
so their induced total-space maps are weak equivalences.
They give singular cohomology isomorphisms, including
the relative disk/sphere groups and their restriction
maps.  Apply the usual mod-2 Thom isomorphism on the
CW model, naturally under the same construction for
\(A_P\).  The rank-zero case is the base pair itself.

Finally, divide the base cuts by \(k^2\) in
\eqref{eq:local-energy-scaling}.  The selected base
label is the root label at \(c_i^s\), so
\eqref{eq:local-thom-dimension} proves the asserted
normal shift.
Iteration is a homeomorphism onto the repeated-loop subspace,
including its oriented phase circles; this identification uses
no covering transfer or division by multiplicity.
\end{proof}

\begin{proof}[Completion of Theorem~\ref{thm:local-control}]
There are only finitely many root levels
\((s\ell_i)^2\), with \(s\mid D_i\).  Around each distinct
level choose a fixed isolated interval with regular
endpoints.  Choose one coarse polygon order \(N_0\)
adequate up to every upper endpoint.

The total crossing pairs in \(P_{N_0}\) have finite
total cohomological rank: their regular sublevels are
compact smooth manifolds with boundary, and the pair
exact sequence applies.  Their relative cohomology
vanishes outside \(0,\ldots,nN_0+1\).
Every selected root label is a direct summand of
such a total crossing by Lemma~\ref{lem:local-lens}.
This includes coincident root levels.

At an arbitrary event, choose the fine polygon order
divisible by \(kN_0\) for every label, where
\(k=m/\gcd(m,D_i)\).  For every even label also impose
divisibility by \(2k'N_0\), where
\[
 k'=\frac{m/2}{\gcd(m/2,D_i)}.
\]
These are finitely many conditions.  After the
normal shift, polygon retractions through the Hilbert
pair identify each finer root crossing with its coarse
pair.  Tight cuts within the fixed isolating intervals
give the same groups by regular-strip deformations.
The finite maximum of their total ranks is therefore
a bound \(B\) for every iterate label.

If the root group in degree \(a\) is nonzero in the
shift calculation, then
\[
 j-I_{i,m}=a-I_{i,s},\qquad 0\le a\le nN_0+1.
\]
For example,
\[
 A=nN_0+1+\max_{i,\ s\mid D_i}I_{i,s}
\]
is a uniform support constant.  Enlarge the constants
if necessary for the fixed computations: on
\(\operatorname{Fix}(\tau)\) start with \(4f_{N/2}\),
apply the same argument to \(m/2\), and divide cuts
by \(4(k')^2\).  The support is centered at
\(I_{i,m/2}\).

The normal shift changes no total rank.  All
degeneracy remains in finitely many root crossings,
so the proof assumes neither nondegeneracy nor
analyticity.  It never divides by \(\Delta_i\);
zero-mean-index primes satisfy the same conclusions.
\end{proof}

The uniform total-rank assertion is for ordinary
local groups.  The half-shift Borel group may have
an infinite polynomial tail.  The Gysin sequence
of Proposition~\ref{prop:gysin} converts the ordinary
bound into a uniform bound in each equivariant
degree and an eventual multiplication isomorphism.
The finite-cut comparison will use these consequences
together with the original, label-preserving
restriction maps established here.
